\origpage{49}{GRAPH MINORS. III}
{\centering
{\Large Graph Minors. III. Planar Tree-Width\par}\medskip

Neil Robertson and P. D. Seymour$^*$\par
\textit{Department of Mathematics, Ohio State University, 231 West 18th Avenue, Columbus, Ohio 43210}\par
\textit{Communicated by the Editors}\par
Received January 14, 1983\par
}
\begin{quote}
The ``tree-width'' of a graph is defined and it is proved that for any fixed planar graph $H$, every planar graph with sufficiently large tree-width has a minor isomorphic to $H$. This result has several applications which are described in other papers in this series.

\end{quote}

\sectionhead{1. \textsc{Introduction}}

All graphs in this paper are finite, and may have loops or multiple edges. $V(G)$ and $E(G)$ denote the sets of vertices and edges of the graph $G$, respectively. A \textit{tree-decomposition} of a graph $G$ is a pair $(T, \mathcal{X})$, where $T$ is a tree and $\mathcal{X} = (X_t : t \in V(T))$ is a family of subsets of $V(G)$, with the following properties:

\hspace*{12pt}(W1) $\bigcup_t (X_t : t \in V(T)) = V(G)$.

\hspace*{12pt}(W2) For every edge $e$ of $G$ there exists $t \in V(T)$ such that $e$ has both ends in $X_t$.

\hspace*{12pt}(W3) For $t, t', t'' \in V(T)$, if $t'$ is on the path of $T$ between $t$ and $t''$ then
\begin{equation*}
X_t \cap X_{t''} \subseteq X_{t'}.
\end{equation*}
The width of the tree-decomposition is
\begin{equation*}
\max_{t \in V(T)} (|X_t| - 1).
\end{equation*}
The graph $G$ has \textit{tree-width} $w$ if $w$ is minimum such that $G$ has a tree-decomposition of width $w$.

Thus, for example, trees and forests have tree-width 1 (or 0 in degenerate%
\origfootnote{Research partially supported by NSF Grant MCS 8103440.}
